% === 親ファイル（ここをコンパイル）========================================
% クラスは自作の styles/thesis.cls を使用（jsbookベース）
\documentclass[twoside,openright]{styles/thesis}

% --- 参考文献（biblatex + biber） -----------------------------------------
% ・backend=biber: コンパイル時に biber を使う
% ・style/citestyle は styles/ 内の自作スタイル（thesis-numeric.bbx）を指定
% ・sorting=none: 引用順
% ・doi, isbn は表示しない。URLは表示する設定
\usepackage[
  backend=biber,
  style=styles/thesis-numeric,
  citestyle=styles/thesis-cite,
  sorting=none, doi=false,
  isbn=false, url=true
]{biblatex}
% 参考文献データベース（.bib）の場所
\addbibresource{backmatter/references.bib}

% --- できるだけ最後に hyperref を読む（リンクの設定）-------------
% dvipdfmx ドライバを使う。リンク色は黒（印刷向け）
\usepackage[
  dvipdfmx,bookmarksnumbered=true,
  colorlinks=true,linkcolor=black,
  citecolor=black,urlcolor=black
]{hyperref}

% --- 共通マクロやコマンドを読み込み ---------------------------------------
% 自作コマンドがあればここにまとめる
% =====================================================================
% preamble/macros.tex
% 目的：
%  - 図表・式などの参照コマンドを統一
% 前提：
%  - styles/thesis-packages.sty 側で graphicx / subfiles を読み込み済み
%  - 本プロジェクトは dvipdfmx（.pdf/.png/.jpg の画像を推奨）
% =====================================================================

% ---------------- 参照系（Fig., Table, Eq. 等を統一） --------------------
% 使い方： \figref{fig:setup}, \tabref{tab:params}, \eqnref{eq:newton}, \chapref{chap:results}, \secref{sec:methods}, \appref{app:hardware}
% 出力例： "Fig. 1", "Table 2", "Eq. (3)", "3章", "2.3節", "付録A"
\newcommand{\figref}[1]{Fig.~\ref{#1}}
\newcommand{\tabref}[1]{Table~\ref{#1}}
\newcommand{\eqnref}[1]{式(\ref{#1})}
\newcommand{\chapref}[1]{\ref{#1}章}
\newcommand{\secref}[1]{\ref{#1}節}
\newcommand{\appref}[1]{付録\ref{#1}}

% ---------------- 画像パス設定（subfiles対応） -------------------------
% \graphicspath + \subfix により、親(main.tex)でも子(chapX.tex)でも
% 「figures/」配下を自動で探索します。
% 例：\includegraphics{chap1/sample_fig.pdf} だけでOK。
% 重要：
%  - 図は figures/chapX/ に置く運用を推奨
\graphicspath{{\subfix{figures/}}}

% ---------------- よく使う書式（必要最小限の便利マクロ） --------------
% ベクトル表記：\vect{x} → 太字 x
\newcommand{\vect}[1]{\mathbf{#1}}

% （任意）よく使う集合：\R, \N
% \newcommand{\R}{\mathbb{R}}
% \newcommand{\N}{\mathbb{N}}

% ---------------- 執筆ルールTIPS（参照ラベル命名規約の提案） ----------
% ラベルは種類ごとに接頭辞を付けると検索しやすい：
%   図     : \label{fig:setup}
%   表     : \label{tab:hyperparam}
%   式     : \label{eq:snr}
%   節/章  : \label{sec:methods}, \label{chap:results}
% 参照例：
%   \figref{fig:setup}に計測システムを示す．\tabref{tab:hyperparam}に計測結果をまとめる．
% =====================================================================

% ========================= ここから本文 ================================
\begin{document}

% --- 前付け ----------------------------
\frontmatter
% 論文表紙、中扉、英文Abstractは教務から配布されるものを差し込む。
\tableofcontents % 目次

% --- 本文：アラビア数字のページ番号をここから振る -------------------------------
\mainmatter
% ========== 第1章 ======================================================
% 章タイトルは親側で管理。本文は子ファイル（chapters/chap1.tex）で編集
\chapter{Introduction} % {}内、章タイトル
\subfile{chapters/chap1}

% ========== 第2章 ======================================================
\chapter{Methods}\label{chap:methods} % labelをつけておくと本文中で参照できる。
\subfile{chapters/chap2}

% ========== 第3章 ======================================================
\chapter{Results}
\subfile{chapters/chap3}

% ========== 第4章 ======================================================
\chapter{Discussion}
\subfile{chapters/chap4}

% ========== 第5章 ======================================================
\chapter{Conclusion}
\subfile{chapters/chap5}


% --- 付録 --------------------------------------------------------------
% \appendix 以降は付録扱い（章番号が A, B, C... になる）
\appendix
\chapter{Additional Figures1}
\subfile{backmatter/appendixA} % ファイルの場所は backmatter/ だが、ページスタイル的に \mainmatter に記述
\chapter{Additional Figures2}
\subfile{backmatter/appendixB}

% --- 後付け ------------------------------------------------------------
\backmatter
% \nocite{*} % 試しに参考文献表示したい時のみコメントアウト外す

% 目次に「参考文献」という項目を手で足す。（biblatex は自動で目次に入れない設定のため）
\addcontentsline{toc}{chapter}{参考文献} 

% 実際の参考文献を出力。見出しの和訳もここで指定
\printbibliography[title={参考文献}]


% 謝辞を差し込み（backmatter/acknowledgment.tex）
% backmatter/acknowledgment.tex
% 謝辞（章番号なし、目次には載せる）
\chapter*{謝辞}
\addcontentsline{toc}{chapter}{謝辞}

本論文の終わりにあたり，本研究に対し終始懇切丁寧なご指導を賜りました 東北大学大学院 工学研究科 田中 真美 教授に厚く御礼申し上げます．

本研究を進めるにあたり，有益なるご教授ならびにご助言を頂きました 東北大学大学院 工学研究科 平田 泰久 教授，山口 健 教授に深く感謝の意を表します．

本研究の遂行に際し，種々ご助力を賜りました 東北大学大学院 工学研究科 奥山 武志 准教授に心より感謝申し上げます．

さらに，本研究に関わる計測実験でご協力頂いた実験参加者の皆様に厚く御礼申し上げます．

最後に公私にわたりお世話になりました 東北大学大学院 工学研究科 ロボティクス専攻 医療福祉工学分野 の皆さまと 東北大学大学院 医工学研究科 医工学専攻 社会医工学講座 医療福祉工学分野 の皆さまに御礼申し上げます．

\bigskip
\noindent
{\raggedleft
2025年10月 \qquad 高成真輝\par}

\end{document}